\section{Постановка задачи}
\label{sec:Chapter1}

\subsection{Объект и цель}

Рассматривается привязной БПЛА, удерживаемый тросом над наземной платформой.
Платформа может быть как стационарной, так и подвижной, поэтому положение
аппарата требуется определять в системе координат, связанной с платформой.
Целью работы является разработка локальной навигационной системы, которая в
реальном времени определяет трёхмерные координаты БПЛА относительно платформы по
ультразвуковым измерениям и сохраняет работоспособность в условиях, где
спутниковая навигация недоступна или ненадёжна.

Формально требуется по набору измерений построить оценку вектора координат
$\mathbf{p}=(x,y,z)$ аппарата в платформенной системе отсчёта, обеспечив заданные
дальность, точность и частоту обновления оценки.

\subsection{Требования к системе}

На основе условий эксплуатации сформулированы следующие требования к
разрабатываемой системе.

\begin{enumerate}
  \item \textbf{Принцип измерения.} Координаты определяются по расстояниям от
        аппарата до нескольких опорных узлов (якорей) с известной геометрией;
        расстояния вычисляются по времени распространения ультразвукового сигнала.
        Это требует общей шкалы времени для передающей и приёмной сторон, то есть
        синхронизации.

  \item \textbf{Рабочая дальность.} Целевая дальность измерений~--- $40$--$50$~м между
        аппаратом и якорями; протокол
        синхронизации проектируется с запасом~--- на дистанции до $100$~м.

  \item \textbf{Точность.} Требуемая погрешность определения координат~--- порядка
        единиц сантиметров, что достаточно для стабилизации аппарата и удержания
        позиции.

  \item \textbf{Частота обновления.} Темп выдачи координат должен соответствовать
        динамике аппарата; для привязного БПЛА с ограниченной скоростью движения
        достаточно частоты порядка $10$~Гц, при этом приёмный тракт не должен
        ограничивать эту величину.

  \item \textbf{Условия эксплуатации.} Система работает в зоне прямой видимости
        якорей, в том числе в стеснённых пространствах с многократными
        переотражениями сигнала, и в непосредственной близости от силового
        оборудования платформы (дизельный электрогенератор, инверторы 400~В),
        создающего сильные электромагнитные помехи. Требуется устойчивость к
        переотражениям и к наводкам по цепям питания и сигнальным линиям.

  \item \textbf{Подвижность платформы.} При возможных наклонах платформы система
        должна позволять пересчёт координат с учётом её ориентации (наличие
        инерциального датчика на платформе).

  \item \textbf{Стоимость и энергопотребление.} Решение должно строиться на
        доступной и недорогой элементной базе с малым энергопотреблением и не
        требовать развёртывания дорогостоящей внешней инфраструктуры.
\end{enumerate}

Перечисленные требования служат критериями при оценке существующих решений
(раздел~\ref{sec:Chapter2}) и при построении предлагаемой системы
(раздел~\ref{sec:Chapter3}).
